% ============================================================================
% LEHRERHINWEISE: HAY (Repaso) + ESTAR
% Klasse 7 | Spanisch | Unidad 2 | DS 02
% ============================================================================

\documentclass[11pt,a4paper]{article}

\usepackage[utf8]{inputenc}
\usepackage[T1]{fontenc}
\usepackage[ngerman]{babel}
\usepackage{geometry}
\usepackage{fancyhdr}
\usepackage{xcolor}
\usepackage{enumitem}
\usepackage{tabularx}
\usepackage{array}
\usepackage{tcolorbox}
\usepackage{booktabs}
\usepackage{amssymb}

% ============================================================================
% LAYOUT
% ============================================================================
\geometry{left=2cm, right=2cm, top=2cm, bottom=2cm}
\definecolor{fach}{HTML}{c41e3a}
\definecolor{grau}{HTML}{666666}
\definecolor{hellgrau}{HTML}{f5f5f5}
\definecolor{basis}{HTML}{2a7a4b}
\definecolor{standard}{HTML}{2c5aa0}
\definecolor{erweiterung}{HTML}{7b1fa2}
\definecolor{spiel}{HTML}{b35c00}

\pagestyle{fancy}
\fancyhf{}
\fancyhead[L]{\small\textcolor{grau}{Spanisch Klasse 7 -- Unidad 2}}
\fancyhead[R]{\small\textcolor{grau}{LH-02 (Lehrerhinweise)}}
\fancyfoot[C]{\small\thepage}
\renewcommand{\headrulewidth}{0.4pt}

% Spanisch-Befehl
\newcommand{\es}[1]{\textbf{#1}}

% ============================================================================
% DOKUMENT
% ============================================================================
\begin{document}

% ----------------------------------------------------------------------------
% KOPFBEREICH
% ----------------------------------------------------------------------------
\begin{center}
{\Large\textcolor{fach}{\textbf{Lehrerhinweise}}}\\[0.3cm]
{\large Hay (Repaso) + Estar}\\[0.2cm]
{\normalsize Doppelstunde (90 Min.) | Qué pasa 1, S. 40--41}
\end{center}

\vspace{0.5cm}
\hrule
\vspace{0.5cm}

% ----------------------------------------------------------------------------
% KURZÜBERSICHT
% ----------------------------------------------------------------------------
\section*{Kurzübersicht}

\begin{tabular}{ll}
\textbf{Thema:} & HAY (Wiederholung + Ergänzung) + ESTAR (neu) \\
\textbf{Lehrwerk:} & Qué pasa 1, Unidad 2, S. 40--41 \\
\textbf{Materialien:} & AB-02, LB-02, SLIDES-02, SPIEL-02 \\
\textbf{Fokus:} & HAY vs ESTAR unterscheiden, Präpositionen \\
\textbf{Voraussetzung:} & Objetos de clase, HAY + un/una/Zahl (aus DS 01) \\
\end{tabular}

\vspace{0.5cm}

% ----------------------------------------------------------------------------
% FEINZIELE
% ----------------------------------------------------------------------------
\section*{Feinziele}

\begin{enumerate}[label=\textbf{FZ\arabic*:}, leftmargin=2.5em]
    \item SuS \textbf{nennen} die Verwendungsregeln von HAY (mit Beispielen) \textit{(AFB I)}
    \item SuS \textbf{bilden} korrekte Sätze mit ESTAR + Ortsangaben \textit{(AFB I)}
    \item SuS \textbf{unterscheiden} HAY und ESTAR in Kontextsätzen \textit{(AFB II)}
    \item SuS \textbf{beschreiben} ihren Klassenraum mit HAY und ESTAR \textit{(AFB II)}
\end{enumerate}

\vspace{0.5cm}

% ----------------------------------------------------------------------------
% STUNDENVERLAUF
% ----------------------------------------------------------------------------
\section*{Stundenverlauf (90 Min.)}

\begin{tabular}{|p{1cm}|p{2.2cm}|p{6.5cm}|p{1cm}|p{2.5cm}|}
\hline
\textbf{Min} & \textbf{Phase} & \textbf{Aktivität} & \textbf{SF} & \textbf{Material} \\
\hline
0--5 & Einstieg & Blitzrunde: \es{¿Qué hay en tu mochila?} (bekannt) & PL & -- \\
\hline
5--15 & Repaso HAY & LB Abschnitt 1: HAY-Regeln (4 Fälle) lesen, Beispiele besprechen & EA/PL & LB-02, SLIDES \\
\hline
15--25 & Übung HAY & AB Aufgabe 1--2: Zuordnung + Sätze bilden & EA & AB-02 \\
\hline
25--40 & Input ESTAR & LB Abschnitt 2: Konjugation + Präpositionen einführen; Vergleich mit QP1 p.41 & PL/EA & LB-02, QP1, SLIDES \\
\hline
40--50 & Übung ESTAR & AB Aufgabe 3: \es{¿Dónde está...?} mit Präpositionen & EA & AB-02 \\
\hline
50--60 & Kontrast & LB Abschnitt 3: HAY vs ESTAR, AB Aufgabe 4: Entscheidung & EA & LB-02, AB-02 \\
\hline
60--70 & Anwendung & AB Aufgabe 5: Klassenraum beschreiben (HAY + ESTAR gemischt) & EA/PA & AB-02 \\
\hline
70--85 & Spiel & \textit{Adivina dónde está} (Ratespiel) & GA & SPIEL-02, Objekte \\
\hline
85--90 & Sicherung & Merksatz formulieren: HAY = Existenz / ESTAR = Ort & PL & -- \\
\hline
\end{tabular}

\vspace{0.5cm}

% ----------------------------------------------------------------------------
% DIFFERENZIERUNG
% ----------------------------------------------------------------------------
\section*{Differenzierung}

\begin{tcolorbox}[colback=white, colframe=basis, title=\textcolor{white}{\textbf{[B] Basis}}]
\begin{itemize}[nosep]
    \item AB Aufgaben 1--3 (nur HAY oder nur ESTAR)
    \item Zuordnung statt freie Sätze
    \item Beim Spiel: Darf Handout benutzen
    \item Kettenübung: Darf Vokabelliste einsehen
\end{itemize}
\end{tcolorbox}

\begin{tcolorbox}[colback=white, colframe=standard, title=\textcolor{white}{\textbf{[S] Standard}}]
\begin{itemize}[nosep]
    \item AB Aufgaben 1--5 vollständig
    \item Freie Sätze mit HAY und ESTAR bilden
    \item Beim Spiel: Alle Präpositionen verwenden
    \item Partnerarbeit ohne Hilfe
\end{itemize}
\end{tcolorbox}

\begin{tcolorbox}[colback=white, colframe=erweiterung, title=\textcolor{white}{\textbf{[E] Erweiterung}}]
\begin{itemize}[nosep]
    \item AB: Alle Aufgaben + Bonus (Mini-Dialog)
    \item Variante: \es{¿Hay...? -- ¿Dónde está?} kombinieren
    \item Expertenrolle: Hilft anderen beim Spiel
    \item Eigenständige Rätsel erfinden
\end{itemize}
\end{tcolorbox}

\vspace{0.5cm}

% ----------------------------------------------------------------------------
% HÄUFIGE FEHLER
% ----------------------------------------------------------------------------
\section*{Häufige Fehler und Reaktionen}

\begin{tabular}{|p{3.8cm}|p{4cm}|p{5.5cm}|}
\hline
\textbf{Fehler} & \textbf{Ursache} & \textbf{Reaktion} \\
\hline
*\es{Hay el libro.} & Deutsch-Transfer & HAY = unbestimmt! Nie mit el/la. \\
\hline
*\es{El libro hay...} & Satzstellung & HAY immer zuerst: \es{Hay un libro.} \\
\hline
*\es{está} für alles & ESTAR übertragen & HAY für Existenz, ESTAR für Ort \\
\hline
\es{estar} ohne Akzent & Akzentregeln unklar & \es{está, estás, estáis} -- Akzent merken! \\
\hline
*\es{en cima de} & Getrenntschreibung & \es{encima de} = ein Wort + de \\
\hline
*\es{a lado de} & al vergessen & \es{al lado de} -- Kontrak\-tion a+el \\
\hline
\end{tabular}

\vspace{0.5cm}

% ----------------------------------------------------------------------------
% SPIELANLEITUNG
% ----------------------------------------------------------------------------
\section*{Spiel: Adivina dónde está}

\begin{tcolorbox}[colback=white, colframe=spiel, title=\textcolor{white}{\textbf{Spielablauf (15 Min.)}}]

\textbf{Material:} 5--6 Klassenzimmerobjekte (estuche, lápiz, goma, libro, regla)

\vspace{0.3cm}

\textbf{Vorbereitung:}
\begin{itemize}[nosep]
    \item SPIEL-02 Handout austeilen (optional)
    \item Präpositionen an Tafel/SLIDES sichtbar
\end{itemize}

\vspace{0.3cm}

\textbf{Ablauf:}
\begin{enumerate}[nosep]
    \item Ein SuS schließt die Augen
    \item Mitschüler verstecken einen Gegenstand (sichtbar, aber ungewöhnlich)
    \item SuS fragt: \es{¿Dónde está el/la...?}
    \item Klasse gibt Hinweise: \es{Está encima de...} / \es{Está al lado de...}
    \item Nach dem Finden: Vollständigen Satz sagen!
\end{enumerate}

\vspace{0.3cm}

\textbf{Variante 2 (schneller):} Lehrer beschreibt Ort $\rightarrow$ SuS erraten Gegenstand.

\vspace{0.3cm}

\textbf{Sprachgerüst:}
\begin{itemize}[nosep]
    \item \es{¿Dónde está el lápiz?}
    \item \es{Está encima de / debajo de / al lado de / delante de / detrás de / entre...}
    \item \es{¡Sí! El lápiz está debajo de la mesa.}
\end{itemize}

\end{tcolorbox}

\vspace{0.5cm}

% ----------------------------------------------------------------------------
% MATERIALCHECKLISTE
% ----------------------------------------------------------------------------
\section*{Materialcheckliste}

\begin{itemize}[nosep]
    \item[$\square$] AB-02-hay-estar (Klassensatz)
    \item[$\square$] LB-02-hay-estar (Klassensatz)
    \item[$\square$] SLIDES-02 (Beamer)
    \item[$\square$] SPIEL-02-adivina (optional, Klassensatz)
    \item[$\square$] Lehrwerk Qué pasa 1 (S. 40--41)
    \item[$\square$] Klassenzimmerobjekte für Spiel (5--6 Stück)
    \item[$\square$] Tafel/Whiteboard für Präpositionen
\end{itemize}

\vspace{0.5cm}

% ----------------------------------------------------------------------------
% TAFELBILD
% ----------------------------------------------------------------------------
\section*{Tafelbild}

\begin{tcolorbox}[colback=white, colframe=grau]
\textbf{HAY vs ESTAR}

\vspace{0.3cm}
\begin{tabular}{p{6.5cm}|p{6.5cm}}
\textbf{HAY = es gibt (Existenz)} & \textbf{ESTAR = sein (Ort/Zustand)} \\
\hline
+ un/una: \es{Hay una mesa.} & estoy, estás, \textbf{está} \\
+ Zahl: \es{Hay tres sillas.} & estamos, estáis, están \\
+ mucho/poco: \es{Hay muchos libros.} & \\
+ ohne Artikel: \es{Hay libros.} & \es{El libro está en la mesa.} \\
& \es{¿Dónde está el lápiz?} \\
\hline
\multicolumn{2}{c}{} \\
\textbf{NICHT:} *hay el libro & \textbf{NICHT:} *el libro hay... \\
\end{tabular}

\vspace{0.3cm}

\textbf{Präpositionen:} en, encima de, debajo de, al lado de, delante de, detrás de, entre
\end{tcolorbox}

\vspace{0.5cm}

% ----------------------------------------------------------------------------
% HAUSAUFGABE
% ----------------------------------------------------------------------------
\section*{Hausaufgabe}

\begin{itemize}[nosep]
    \item Präpositionen lernen (en, encima de, debajo de, al lado de, delante de, detrás de)
    \item Lehrwerk S. 41, Aufgabe 3
    \item Optional: LB-02 nochmal durchlesen
    \item [E]: 5 Sätze über eigenes Zimmer mit HAY + ESTAR
\end{itemize}

\vspace{0.5cm}

% ----------------------------------------------------------------------------
% AB-LÖSUNGEN
% ----------------------------------------------------------------------------
\section*{AB-Lösungen}

\textbf{Aufg. 1 -- Zuordnung:}
\begin{itemize}[nosep]
    \item hay + un/una $\rightarrow$ Hay una mesa.
    \item hay + número $\rightarrow$ Hay tres sillas.
    \item hay + mucho/poco $\rightarrow$ Hay muchos libros.
    \item hay + sin artículo $\rightarrow$ Hay libros.
\end{itemize}

\vspace{0.3cm}

\textbf{Aufg. 2 -- Sätze mit hay:}
\begin{enumerate}[label=\alph*), nosep]
    \item hay una pizarra
    \item hay veinticinco sillas
    \item hay muchos bolígrafos
    \item hay libros
    \item hay pocos cuadernos
    \item hay un ordenador
\end{enumerate}

\vspace{0.3cm}

\textbf{Aufg. 3 -- estar + Präposition:}
\begin{enumerate}[label=\alph*), nosep]
    \item está
    \item está encima de / está en
    \item está debajo de / está al lado de
    \item está delante de
    \item están en
    \item estoy al lado de / estoy delante de
\end{enumerate}

\vspace{0.3cm}

\textbf{Aufg. 4 -- hay o estar:}
\begin{enumerate}[label=\alph*), nosep]
    \item \textbf{Hay} un libro... $\rightarrow$ hay (unbestimmt: \textit{un})
    \item ¿Dónde \textbf{está} la goma? $\rightarrow$ estar: está (bestimmt: \textit{la})
    \item En la clase \textbf{hay} muchas... $\rightarrow$ hay (Existenz)
    \item El cuaderno \textbf{está}... $\rightarrow$ estar: está (bestimmt: \textit{el})
    \item ¿\textbf{Hay} lápices...? $\rightarrow$ hay (ohne Artikel)
    \item Mi mochila \textbf{está}... $\rightarrow$ estar: está (bestimmt: \textit{mi})
\end{enumerate}

\vspace{0.3cm}

\textbf{Aufg. 5 -- Beschreibung:} Individuelle Sätze. Kriterien:
\begin{itemize}[nosep]
    \item 4× hay (mit verschiedenen Strukturen)
    \item 4× estar (mit verschiedenen Präpositionen)
    \item Korrekte Artikel (el/la/un/una)
\end{itemize}

\vspace{0.3cm}

\textbf{Bonus -- Mini-Dialog:} Individuelle Dialoge. Muster:
\begin{itemize}[nosep]
    \item A: \es{¿Hay un bolígrafo?}
    \item B: \es{Sí, hay uno.}
    \item A: \es{¿Dónde está?}
    \item B: \es{Está encima de la mesa.}
\end{itemize}

\end{document}
